\chapter{Atomic, diffuse, and singular measures}\label{sec:examples-measures}
These examples make the distinctions between mass, support, density, and weak convergence explicit. The conventions are Definitions \ref{def:radon-measure}, \ref{def:weak-measure}, and \ref{def:hausdorff-measure}. The foundational source is \cite{Sim18}, Chapter 1, Sections 1--5. The Cantor constructions below are proved directly; they illustrate the covering definitions rather than invoke a structure theorem for fractals.

\begin{example}[A measure can be atomic and have an interval as support]\label{ex:dense-atomic-measure}
Enumerate the rational points of $[0,1]$ without repetitions as $q_1,q_2,\ldots$, and set
\[
 \mu=\sum_{j=1}^{\infty}2^{-j}\delta_{q_j}.
\]
This is a Radon probability measure on $\mathbb R$, concentrated on a countable set, with support $[0,1]$. At each $q_j$, its one-dimensional density is infinite. In contrast, counting measure on $\mathbb Q\cap[0,1]$ is not locally finite.
\end{example}
\begin{proof}
Countable additivity follows by interchanging nonnegative sums. Finite partial sums are finite Radon measures, and their omitted tails have total mass tending to zero, which proves inner and outer regularity of $\mu$ by approximation. Every neighborhood of a point of $[0,1]$ contains some $q_j$ and hence positive mass; outside that interval there are neighborhoods of zero mass. At $q_j$, $\mu(B_r(q_j))/(2r)\geq2^{-j}/(2r)\to\infty$. Counting measure assigns infinite mass to the compact interval $[0,1]$. Thus support alone records neither atomicity nor local finiteness.
\end{proof}

\begin{example}[Concentration of absolutely continuous measures]\label{ex:approximate-identity-dirac}
Let $\rho\geq0$ belong to $L^1(\mathbb R^n)$ with $\int\rho=1$, fix $a\in\mathbb R^n$, and let $\varepsilon\downarrow0$. Then
\[
 d\mu_\varepsilon(x)=\varepsilon^{-n}\rho((x-a)/\varepsilon)\,dx
 \quad\Longrightarrow\quad \mu_\varepsilon\rightharpoonup\delta_a.
\]
Every approximating measure is absolutely continuous with respect to Lebesgue measure, while the limit is singular and atomic.
\end{example}
\begin{proof}
For $\varphi\in C_c(\mathbb R^n)$, changing variables gives $\int\varphi\,d\mu_\varepsilon=\int\varphi(a+\varepsilon y)\rho(y)\,dy$. The integrand converges pointwise to $\varphi(a)\rho(y)$ and is bounded by $\|\varphi\|_\infty\rho(y)$. Dominated convergence proves the claim. Local finiteness and regularity of the approximating measures follow from integration of a locally integrable density against Lebesgue measure.
\end{proof}

\begin{example}[Mass escapes through infinity or a domain boundary]\label{ex:escaping-mass}
In $\mathbb R$, $\delta_j\rightharpoonup0$ although each measure has total mass one. In the open domain $U=(0,1)$, $\delta_{1/j}\rightharpoonup0$ for $j\geq2$, whereas the same measures viewed on $\mathbb R$ converge to $\delta_0$.
\end{example}
\begin{proof}
A compactly supported continuous test vanishes at $j$ for all sufficiently large $j$. A compact subset of $(0,1)$ has positive distance from zero, so a test supported there also eventually vanishes at $1/j$. On $\mathbb R$, continuity instead gives $\varphi(1/j)\to\varphi(0)$. The choice of the ambient open set is part of Definition \ref{def:weak-measure}; preservation of total mass does not follow from that topology.
\end{proof}

\begin{example}[Weak convergence does not imply convergence on every Borel set]\label{ex:weak-borel-discontinuity}
The sequence $\delta_{1/j}\rightharpoonup\delta_0$ on $\mathbb R$ has
\[
 \delta_{1/j}(\{0\})=0,\qquad\delta_0(\{0\})=1,
 \qquad \delta_{1/j}((0,1))=1\quad(j\geq2),\qquad\delta_0((0,1))=0.
\]
Thus both a closed set and an open set can have discontinuous masses. The open-set lower-semicontinuity inequality in Proposition \ref{prop:measure-lsc} remains valid.
\end{example}
\begin{proof}
Evaluation of a Dirac measure at a set is its indicator function at the atom. The four values follow immediately. Here zero lies on the boundary of each of the sets at issue and carries positive limiting mass; indicators of these sets are not continuous tests at that point.
\end{proof}

\begin{example}[Oscillating densities can average without converging in norm]\label{ex:oscillating-densities}
On $[0,1]$, put $f_j(x)=1+\sin(2\pi jx)$ and $d\mu_j=f_j\,dx$. Then $\mu_j$ are probability measures and converge weakly on $\mathbb R$ to $\mathcal L^1\llcorner[0,1]$, but
\[
 \int_0^1|f_j-1|\,dx=\frac2\pi
\]
for every positive integer $j$. In particular weak convergence here does not imply convergence of densities in $L^1$.
\end{example}
\begin{proof}
Nonnegativity and mass one follow by direct integration. For continuous $\varphi$ on $[0,1]$, partition into the $j$ intervals $I_k=[k/j,(k+1)/j]$. The sine has zero integral on each. Subtracting $\varphi(k/j)$ on $I_k$ bounds the oscillating integral in absolute value by the modulus of continuity of $\varphi$ at $1/j$, which tends to zero. The displayed norm is $j$ equal copies of one period, rescaled to total length one; the integral of $|\sin|$ over a period is four.
\end{proof}

\begin{definition}[The middle-thirds Cantor set and its probability measure]\label{def:cantor-model}
Let $C_k$ be the union of the $2^k$ closed intervals of length $3^{-k}$ obtained by retaining the left and right thirds at each of $k$ steps, and put $C=\bigcap_{k\geq0}C_k$. For $t\in[0,1)$ write its binary digits as $b_j(t)\in\{0,1\}$, choosing the expansion eventually zero at dyadic rationals. Define
\[
 T(t)=\sum_{j=1}^{\infty}\frac{2b_j(t)}{3^j},
 \qquad \mu_C=T_\#(\mathcal L^1\llcorner[0,1)).
\]
The pushforward means $\mu_C(A)=\mathcal L^1(T^{-1}(A))$ for Borel sets $A$. Each finite digit function is measurable, so the limit $T$ is measurable. This constructs the measure without presupposing a measure on an infinite product space.
\end{definition}

\begin{proposition}[Cylinder masses and singular continuity]\label{prop:cantor-cylinder-mass}
The set $C$ is compact and has Lebesgue measure zero. The measure $\mu_C$ is a Radon probability measure supported on $C$, and every level-$k$ retained interval $I$ has mass $2^{-k}$. It has no atoms. Thus a measure may be singular with respect to Lebesgue measure without being atomic.
\end{proposition}
\begin{proof}
The sets $C_k$ are nested compact sets with total length $(2/3)^k$, so their intersection is compact and null. The ternary series defining $T$ belongs to every $C_k$. Specifying a word of $k$ binary digits selects a half-open dyadic interval of Lebesgue measure $2^{-k}$ and the corresponding retained ternary interval; endpoint conventions affect at most a countable null set. This proves the cylinder mass. Every point of $C$ lies in a retained interval at every level; its singleton mass is therefore at most $2^{-k}$ for all $k$. These same shrinking intervals show that every neighborhood of every point of $C$ has positive mass, so the support is exactly $C$. A finite Borel measure on $\mathbb R$ is regular (\cite{Sim18}, Chapter 1, Section 1); hence this probability measure is Radon. Alternatively, regularity follows by compact cylinder approximations and finite total mass.
\end{proof}

\begin{proposition}[Power-law bounds for the Cantor measure]\label{prop:cantor-ball-bounds}
Write $s=\log2/\log3$. For every $x\in C$ and $0<r<1$,
\[
 \frac18r^s\leq\mu_C(B_r(x))\leq8r^s.
\]
In particular $\mu_C(B_r(x))/(2r)\to\infty$ at every $x\in C$, despite the absence of atoms.
\end{proposition}
\begin{proof}
Choose $k\geq1$ with $3^{-k}\leq r<3^{-k+1}$. The level-$k+2$ interval containing $x$ has diameter strictly less than $r$ and hence lies in $B_r(x)$. Its mass is $2^{-k-2}$. Since $r^s<2\cdot2^{-k}$, this is at least $r^s/8$. Distinct level-$k$ intervals have length $L=3^{-k}$ and gaps at least $L$. An interval of length $2r<6L$ meets at most eight of them. Their total mass is at most $8\cdot2^{-k}\leq8r^s$. Finally $s<1$, so the lower bound divided by $2r$ diverges. The constants are deliberately nonsharp.
\end{proof}

\begin{definition}[Hausdorff dimension via diameter covers]\label{def:hausdorff-dimension-examples}
For a metric subset $E$ and $q>0$, let
\[
 \mathcal H^q_{\rm diam}(E)=\lim_{\delta\downarrow0}
 \inf\left\{\sum_i(\operatorname{diam}E_i)^q:
 E\subset\bigcup_iE_i,\ \operatorname{diam}E_i<\delta\right\}.
\]
Define $\dim_H E=\inf\{q>0:\mathcal H^q_{\rm diam}(E)=0\}$. For integer $q=m\geq1$, this measure differs from Definition \ref{def:hausdorff-measure} by the positive constant $2^m/\omega_m$; this constant does not affect dimension. The covering construction for real exponents is the one in \cite{Sim18}, Chapter 1, Section 2.
\end{definition}

\begin{proposition}[The Cantor set has a noninteger dimension]\label{prop:cantor-dimension}
For $s=\log2/\log3$,
\[
 0<\mathcal H^s_{\rm diam}(C)<\infty,\qquad \dim_H C=s.
\]
Thus zero length does not imply zero Hausdorff dimension.
\end{proposition}
\begin{proof}
The level-$k$ cover has sum $2^k(3^{-k})^s=1$, and for $q>s$ its $q$-sum tends to zero. For the lower bound take any sufficiently fine countable cover. A set of diameter $d>0$ meeting $C$ is contained, on $C$, in the ball of radius $2d$ about any chosen point of its intersection. Proposition \ref{prop:cantor-ball-bounds} bounds its $\mu_C$ outer measure by $8(2d)^s$. Sets of diameter zero carry no mass. Countable subadditivity gives $1\leq8\cdot2^s\sum_i d_i^s$, hence a uniform positive lower bound. If $0<q<s$ and $d_i<\delta$, then $\sum_i d_i^q\geq\delta^{q-s}\sum_i d_i^s$, which tends to infinity as $\delta\downarrow0$. This also verifies the dimension assertion directly.
\end{proof}

\begin{example}[The Cantor function has a singular distributional derivative]\label{ex:cantor-function-bv}
Let $F(x)=\mu_C(( -\infty,x])$. Then $F$ is continuous, nondecreasing, equals zero for $x\leq0$ and one for $x\geq1$, and is constant on every component of $\mathbb R\setminus C$. It satisfies
\[
 F'=0\quad\mathcal L^1\text{-almost everywhere},\qquad DF=\mu_C
\]
as distributions. In particular it is locally BV but is not absolutely continuous on $[0,1]$.
\end{example}
\begin{proof}
A distribution function jumps precisely by its singleton mass, so Proposition \ref{prop:cantor-cylinder-mass} gives continuity. An interval disjoint from the support has zero mass, giving constancy there and zero classical derivative off the null set $C$. For $\varphi\in C_c^\infty(\mathbb R)$, Fubini yields
\[
 -\int F(x)\varphi'(x)\,dx
 =-\int_C\int_t^\infty\varphi'(x)\,dx\,d\mu_C(t)
 =\int\varphi(t)\,d\mu_C(t).
\]
The finite measure derivative proves local BV. Absolute continuity on $[0,1]$ would imply $F(1)-F(0)=\int_0^1F'=0$, contradicting the endpoint values. This is an explicit singular continuous part of a BV derivative, distinct from the jump part of a step function.
\end{proof}


\begin{example}[Diffuse approximations to a singular continuous measure]\label{ex:cantor-density-approximation}
Let $C_k$ be the level-$k$ Cantor approximation and define
\[
 d\nu_k=(3/2)^k\mathbf1_{C_k}\,dx.
\]
Then $\nu_k$ is an absolutely continuous probability measure and $\nu_k\rightharpoonup\mu_C$. Each retained interval has the same mass under $\nu_k$ and $\mu_C$, although their behavior inside it is different.
\end{example}
\begin{proof}
The total length of $C_k$ is $(2/3)^k$. Multiplying by $(3/2)^k$ gives mass one, and each interval has mass $2^{-k}$. For a continuous test $\varphi$ on $[0,1]$, select one point $x_I$ in each level-$k$ interval. The difference between either measure's integral and $\sum_I2^{-k}\varphi(x_I)$ is bounded by the modulus of continuity of $\varphi$ at $3^{-k}$. These two errors tend to zero, proving weak convergence. Thus concentration can produce a singular continuous limit as well as the Dirac limit in Example \ref{ex:approximate-identity-dirac}.
\end{proof}

\begin{example}[Three mutually singular components in one measure]\label{ex:three-measure-components}
The finite measure
\[
 \eta=\mathcal L^1\llcorner[0,1]+\mu_C+\delta_0
\]
has an absolutely continuous part, a singular continuous part, and an atomic part. These three summands are mutually singular even though their supports overlap. Its distribution function on $\mathbb R$ is the sum of the distribution functions of the three summands.
\end{example}
\begin{proof}
The three measures are respectively concentrated on the disjoint Borel sets $[0,1]\setminus C$, $C\setminus\{0\}$, and $\{0\}$. Removing $C$ does not change Lebesgue measure, and removing zero does not change $\mu_C$ because it has no atoms. Each summand is therefore singular to the others, while the first two supports are $[0,1]$ and $C$. Additivity gives the distribution-function assertion. This explicit decomposition does not require the general Lebesgue decomposition theorem.
\end{proof}
